% A rigid gate sweep through one isolated charge transition.
% All coordinates are schematic; the three settings are consistent
 % with the zero-temperature off--on--off current in the lower panel.
\begingroup
\definecolor{transportA}{HTML}{0072B2}
\definecolor{transportB}{HTML}{D55E00}
\definecolor{transportC}{HTML}{009E73}
% Logistic step, written using only nonpositive exponential arguments.
\pgfmathdeclarefunction{transportStep}{1}{%
 \pgfmathparse{0.5*(1+sign(#1)*(1-exp(-abs(#1)))/(1+exp(-abs(#1))))}%
}
\begin{tikzpicture}[
 x=1mm,y=1mm,>=stealth,line cap=round,line join=round,
 every node/.style={font=\small,inner sep=0pt}
]
 \path[use as bounding box] (0,8) rectangle (150,123);

 \begin{scope}
  \clip (17,56) rectangle (133,116);
  \fill[white] (17,56) rectangle (133,116);
  % The five thresholds have the same color at a given gate setting.
  \foreach \thresholdX/\thresholdY in {30/63,53/68,76/86,99/104,122/109}{
   \fill[transportA] (\thresholdX,{\thresholdY+6}) circle[radius=0.65mm];
   \fill[transportB] (\thresholdX,\thresholdY) circle[radius=0.65mm];
   \fill[transportC] (\thresholdX,{\thresholdY-6}) circle[radius=0.65mm];
  }
  % Real translucent masks are drawn AFTER the colored points.
  % The unmasked white strip is the window 83 < mu < 89.
  \fill[black!15,fill opacity=0.45,draw=none]
       (17,56) rectangle (133,83);
  \fill[black!15,fill opacity=0.45,draw=none]
       (17,89) rectangle (133,116);
 \end{scope}

 % Arrows between the target thresholds indicate increasing gate voltage.
 \draw[->,black!55,line width=0.45pt] (76,90.9)--(76,87.1);
 \draw[->,black!55,line width=0.45pt] (76,84.9)--(76,81.1);

 % No horizontal boundary strokes: shading alone delimits the window.
 \draw[->,line width=0.45pt] (17,56)--(136,56);
 \draw[->,line width=0.45pt] (17,56)--(17,119);
 \node[anchor=west] at (0,116) {(a)};
 \node[anchor=east] at (13,119) {$\mu_n$};
 \node[anchor=west] at (140,56) {$n$};
 \node[anchor=west] at (136,89) {$\mu_{\mathrm S}$};
 \node[anchor=west] at (136,83) {$\mu_{\mathrm D}$};
 \foreach \tickX/\tickLabel in {30/{N-2},53/{N-1},76/{N},99/{N+1},122/{N+2}}{
  \draw[line width=0.4pt] (\tickX,56)--(\tickX,54.8);
  \node at (\tickX,51.5) {$\tickLabel$};
 }

 % The plotted linear response is mu_N=86-3*(x-75)/17.
 % Thus mu_N crosses source 89 at x=58 and drain 83 at x=92.
 % The snapshots x=41,75,109 correspond to energy shifts +6,0,-6.
 \node[anchor=west] at (0,40) {(b)};
 \draw[->,black!70,line width=0.45pt] (17,13)--(136,13);
 \draw[->,black!70,line width=0.45pt] (17,13)--(17,42);
 \node[anchor=east] at (13,42) {$I$};
 \node[anchor=east] at (13,13) {$0$};
 \node[anchor=west] at (140,13) {$V_g$};
 \draw[black!85,line width=0.8pt]
      (21,13)--(58,13)--(58,35)--(92,35)--(92,13)--(131,13);
 % Finite-T current is the Fermi-function difference f_S-f_D.
 % In the rescaled gate coordinate, the thermal width is 1.7:
 % f_S-f_D = logistic((x-58)/1.7)-logistic((x-92)/1.7).
 % This is k_B T_e/(e V_b)=1.7/34=0.05, with the same unnormalized
 % current scale and tunnel rates as the zero-temperature trace.
 \draw[black!85,line width=0.85pt,dash pattern=on 2.3pt off 1.7pt]
      plot[domain=21:131,samples=221]
      (\x,{13+22*(transportStep((\x-58)/1.7)-transportStep((\x-92)/1.7))});
 \draw[black!85,line width=0.8pt] (105,39)--(113,39);
 \node[anchor=west] at (116,39) {$T_e=0$};
 \draw[black!85,line width=0.85pt,dash pattern=on 2.3pt off 1.7pt]
      (105,33)--(113,33);
 \node[anchor=west] at (116,33) {$T_e>0$};
\end{tikzpicture}
\endgroup
